%-------------------------------------------------------------------------------
% Abstract
%-------------------------------------------------------------------------------

\chapter*{Abstract}
\addcontentsline{toc}{chapter}{Abstract}

Automated Algorithm Design (AAD) aims to reduce the manual effort required to develop optimization algorithms by automating the design of the search procedure itself. Large Language Models (LLMs) have recently made it possible to generate executable optimization algorithms from natural-language descriptions. This raises the question of whether such models can generate useful heuristics for continuous black-box optimization when fitness evaluations are affected by stochastic noise.

This thesis investigates LLM-based synthesis of continuous optimization heuristics under deterministic and noisy conditions using the Large Language Model Evolutionary Algorithm (LLaMEA). Candidate Python optimizers are generated and evaluated on continuous Black-Box Optimization Benchmark (BBOB) functions with controlled heteroscedastic fitness noise. Two model scales and four prompt configurations are considered during synthesis. The resulting heuristics are then independently evaluated alongside CMA-ES, Differential Evolution, and Particle Swarm Optimization on representative unimodal, ill-conditioned, and multimodal problems in 2D, 5D, and 10D, using four noise levels ranging from deterministic to increasingly noisy evaluations.

Performance is evaluated through optimality gaps, success rates, convergence trajectories, target-attainment ECDFs, and anytime-performance measures over repeated trials. The experiments focus on how model scale, prompt configuration, noise level, problem structure, and dimensionality affect the generated heuristics. The results are used to compare their performance with established optimization methods and to identify the conditions under which LLM-generated algorithms succeed, degrade, or fail.
